\section{Korzyści zastosowania metod zespołowego uczenia maszynowego}
Jak wspomniano w podanym artykule \cite{metoda13},
\textbf{uczenie zespołowe} to metoda uczenia maszynowego, która polega
na łączeniu wyników wielu modeli w celu poprawy dokładności predykcji. Główne klasy metod uczenia zespołowego obejmują
\textbf{agregację metodą bootstrap (bagging)}, \textbf{uczenie wzmacniające (boosting)} oraz \textbf{warstwowe łączenie modeli (stacking)}. Bardziej
szczegółowy opis wspominanych powyżej metod zostanie omówiony w
następnym podrozdziale.

W porównaniu do pojedynczych modeli, metody zespołowe osiągają lepsze
wyniki zarówno w metrykach klasyfikacyjnych, jak i ogólnej dokładności.
Ponadto, odpowiednio zaprojektowane metody ensemble mogą redukować
problem nadmiernego dopasowania (overfitting). Chciałbym wyróżnić
następujące \textbf{zalety} stosowania metod zespołowego
uczenia maszynowego:

\begin{itemize}
	\item
	      \textbf{Zwiększona odporność na błędy}: uczenie zespołowe zmniejsza
	      ryzyko błędów spowodowanych nadmierną wariancją lub stronniczością
	      pojedynczych modeli, co zapewnia bardziej stabilne i wiarygodne
	      wyniki.
	\item
	      \textbf{Integracja różnorodnych modeli}: Możliwość łączenia różnych
	      typów modeli pozwala na wykorzystanie ich indywidualnych mocnych
	      stron, co zwiększa skuteczność w zadaniach
	      klasyfikacyjnych i regresyjnych.
	\item
	      \textbf{Skuteczność w rozwiązywaniu problemów nierównowagi klas}:
	      uczenie zespołowe wykazuje lepsze wyniki w zadaniach z nierównomiernym
	      rozkładem klas, zwłaszcza gdy jest stosowany w połączeniu z metodami
	      wzbogacania danych, takimi jak \textbf{SMOTE}.
\end{itemize}
